\documentclass[10pt,a4paper]{amsart}
\usepackage{amsmath,amssymb,amsthm,graphicx}
\usepackage{verbatim}
\newenvironment{reference}{\comment}{\endcomment}
%\newenvironment{reference}{}{}
\newenvironment{slogan}{\comment}{\endcomment}
\newenvironment{history}{\comment}{\endcomment}

% For commutative diagrams we use Xy-pic
\usepackage[all]{xy}

% We use 2cell for 2-commutative diagrams.
\xyoption{2cell}
% Theorem environments.
%
\theoremstyle{plain}
\newtheorem{theorem}[subsection]{Theorem}
\newtheorem{proposition}[subsection]{Proposition}
\newtheorem{lemma}[subsection]{Lemma}

\theoremstyle{definition}
\newtheorem{definition}[subsection]{Definition}
\newtheorem{example}[subsection]{Example}
\newtheorem{exercise}[subsection]{Exercise}
\newtheorem{situation}[subsection]{Situation}

\theoremstyle{remark}
\newtheorem{remark}[subsection]{Remark}
\newtheorem{remarks}[subsection]{Remarks}

\numberwithin{equation}{subsection}

% Macros
%
\def\lim{\mathop{\mathrm{lim}}\nolimits}
\def\colim{\mathop{\mathrm{colim}}\nolimits}
\def\Spec{\mathop{\mathrm{Spec}}}
\def\Hom{\mathop{\mathrm{Hom}}\nolimits}
\def\Ext{\mathop{\mathrm{Ext}}\nolimits}
\def\SheafHom{\mathop{\mathcal{H}\!\mathit{om}}\nolimits}
\def\SheafExt{\mathop{\mathcal{E}\!\mathit{xt}}\nolimits}
\def\Sch{\mathit{Sch}}
\def\Mor{\mathop{\mathrm{Mor}}\nolimits}
\def\Ob{\mathop{\mathrm{Ob}}\nolimits}
\def\Sh{\mathop{\mathit{Sh}}\nolimits}
\def\NL{\mathop{N\!L}\nolimits}
\def\CH{\mathop{\mathrm{CH}}\nolimits}
\def\proetale{{pro\text{-}\acute{e}tale}}
\def\etale{{\acute{e}tale}}
\def\QCoh{\mathit{QCoh}}
\def\Ker{\mathop{\mathrm{Ker}}}
\def\Im{\mathop{\mathrm{Im}}}
\def\Coker{\mathop{\mathrm{Coker}}}
\def\Coim{\mathop{\mathrm{Coim}}}

% Boxtimes
%
\DeclareMathSymbol{\boxtimes}{\mathbin}{AMSa}{"02}

%
% Macros for moduli stacks/spaces
%
\def\QCohstack{\mathcal{QC}\!\mathit{oh}}
\def\Cohstack{\mathcal{C}\!\mathit{oh}}
\def\Spacesstack{\mathcal{S}\!\mathit{paces}}
\def\Quotfunctor{\mathrm{Quot}}
\def\Hilbfunctor{\mathrm{Hilb}}
\def\Curvesstack{\mathcal{C}\!\mathit{urves}}
\def\Polarizedstack{\mathcal{P}\!\mathit{olarized}}
\def\Complexesstack{\mathcal{C}\!\mathit{omplexes}}
% \Pic is the operator that assigns to X its picard group, usage \Pic(X)
% \Picardstack_{X/B} denotes the Picard stack of X over B
% \Picardfunctor_{X/B} denotes the Picard functor of X over B
\def\Pic{\mathop{\mathrm{Pic}}\nolimits}
\def\Picardstack{\mathcal{P}\!\mathit{ic}}
\def\Picardfunctor{\mathrm{Pic}}
\def\Deformationcategory{\mathcal{D}\!\mathit{ef}}
\usepackage{cleveref}
\crefname{originalreference}{original reference}{original references}
\Crefname{originalreference}{Original reference}{Original references}
\makeatletter
\newlabel{situation-quasi-compact-objects}{{0937}{source}{}{sourcepointer-1}{}}
\newlabel{situation-quasi-compact-objects@cref}{{[originalreference][0][]0937}{[1][0][]source}{}{}{}}
\makeatother
\crefname{originalreference}{original reference}{original references}
\Crefname{originalreference}{Original reference}{Original references}
\makeatletter
\newlabel{item-enough}{{0938}{source}{}{sourcepointer-2}{}}
\newlabel{item-enough@cref}{{[originalreference][0][]0938}{[1][0][]source}{}{}{}}
\makeatother
\crefname{originalreference}{original reference}{original references}
\Crefname{originalreference}{Original reference}{Original references}
\makeatletter
\newlabel{item-enough-qc}{{0939}{source}{}{sourcepointer-3}{}}
\newlabel{item-enough-qc@cref}{{[originalreference][0][]0939}{[1][0][]source}{}{}{}}
\makeatother
\crefname{originalequation}{original equation}{original equations}
\Crefname{originalequation}{Original equation}{Original equations}
\makeatletter
\newlabel{equation-towards-constructible}{{093D}{source}{}{sourcepointer-4}{}}
\newlabel{equation-towards-constructible@cref}{{[originalequation][0][]093D}{[1][0][]source}{}{}{}}
\makeatother
\crefname{originallemma}{original lemma}{original lemmas}
\Crefname{originallemma}{Original lemma}{Original lemmas}
\makeatletter
\newlabel{lemma-silly-quasi-compact}{{0G1W}{source}{}{sourcepointer-5}{}}
\newlabel{lemma-silly-quasi-compact@cref}{{[originallemma][0][]0G1W}{[1][0][]source}{}{}{}}
\makeatother
\title{Cokernels of maps between the standard presentations are again standard presentations}
\author{The Stacks Project Authors}
\date{stacks-project:sites-modules.tex:local-snapshot-20261012; GNU FDL 1.2 or later}
\begin{document}
\maketitle
\noindent\textit{Editorial note.} All delivered mathematical text is the author's own, copied verbatim from the pinned Stacks Project chapter file sites-modules.tex; no statement or proof step was rewritten, reordered or removed.\par\medskip
\noindent\textit{Editorial note.} The counted span cites five external items (a situation, several lemmas and an equation), all recorded in assembly.external\_references with canonical tags as navigation pointers.\par\medskip
\noindent\textit{Allowed standard prerequisites (editorial).}\begin{itemize}
\item Finite subfamilies of a basis covering of a quasi-compact object. Stacks Project, Modules on Sites, Lemma lemma-silly-quasi-compact (tag 09B0). Produces the finite subset K' used to exhibit the presentation.
\item The quasi-compact objects situation and its presentation equation. Stacks Project, Modules on Sites, Situation situation-quasi-compact-objects (tag 09AZ). Supplies the basis, the enough/ quasi-compact hypotheses and the standard presentation whose closure is proved.
\end{itemize}\medskip
\ifdefined\hypertarget\hypertarget{sourcepointer-1}{}\fi
\noindent\textit{Source reference.} Original reference 0937: sources/stacks-chow/tags.txt, source lines 11505--11505. Cross-chapter standard prerequisite cited by the author inside the counted span; it is used as a pointer and its own proof is not part of this item.\par\medskip
\ifdefined\hypertarget\hypertarget{sourcepointer-2}{}\fi
\noindent\textit{Source reference.} Original reference 0938: sources/stacks-chow/tags.txt, source lines 11506--11506. Cross-chapter standard prerequisite cited by the author inside the counted span; it is used as a pointer and its own proof is not part of this item.\par\medskip
\ifdefined\hypertarget\hypertarget{sourcepointer-3}{}\fi
\noindent\textit{Source reference.} Original reference 0939: sources/stacks-chow/tags.txt, source lines 11507--11507. Cross-chapter standard prerequisite cited by the author inside the counted span; it is used as a pointer and its own proof is not part of this item.\par\medskip
\ifdefined\hypertarget\hypertarget{sourcepointer-4}{}\fi
\noindent\textit{Source reference.} Original equation 093D: sources/stacks-chow/tags.txt, source lines 11511--11511. Cross-chapter standard prerequisite cited by the author inside the counted span; it is used as a pointer and its own proof is not part of this item.\par\medskip
\ifdefined\hypertarget\hypertarget{sourcepointer-5}{}\fi
\noindent\textit{Source reference.} Original lemma 0G1W: sources/stacks-chow/tags.txt, source lines 20036--20036. Cross-chapter standard prerequisite cited by the author inside the counted span; it is used as a pointer and its own proof is not part of this item.\par\medskip
\begin{lemma}
\label{lemma-cokernel-map-towards-constructibles}
In Situation \ref{situation-quasi-compact-objects} assume
(\ref{item-enough}) and (\ref{item-enough-qc}) hold.
Let $\mathcal{O}$ be a sheaf of rings.
Then a cokernel of a map between modules as in
(\ref{equation-towards-constructible}) is another module as
in (\ref{equation-towards-constructible}).
\end{lemma}

\begin{proof}
Let $\mathcal{F} = \Coker(\bigoplus j_{V_j!}\mathcal{O}_{V_j} \to
\bigoplus j_{U_i!}\mathcal{O}_{U_i})$
as in (\ref{equation-towards-constructible}). It suffices to show
that the cokernel of a map $\varphi : j_{W!}\mathcal{O}_W \to \mathcal{F}$
with $W \in \mathcal{B}$ is another module of the same type.
The map $\varphi$ corresponds to $s \in \mathcal{F}(W)$.
Since $\bigoplus j_{U_i!}\mathcal{O}_{U_i} \to \mathcal{F}$ is
surjective, by (\ref{item-enough}) we may choose a covering
$\{W_k \to W\}_{k \in K}$ with $W_k \in \mathcal{B}$
such that $s|_{W_k}$ is the image of some section
$s_k$ of $\bigoplus j_{U_i!}\mathcal{O}_{U_i})$.
By (\ref{item-enough-qc}) the object $W$ is quasi-compact.
By Lemma \ref{lemma-silly-quasi-compact}
there is a finite subset $K' \subset K$ such that
$\bigoplus_{k \in K'} j_{W_k!}\mathcal{O}_{W_k} \to j_{W!}\mathcal{O}_W$
is surjective. We conclude that $\Coker(\varphi)$ is equal to
$$
\Coker\left(
\bigoplus\nolimits_{k \in K'} j_{W_k!}\mathcal{O}_{W_k} \oplus
\bigoplus j_{V_j!}\mathcal{O}_{V_j}
\longrightarrow
\bigoplus j_{U_i!}\mathcal{O}_{U_i}
\right)
$$
where the map $\bigoplus_{k \in K'} j_{W_k!}\mathcal{O}_{W_k}
\to \bigoplus j_{U_i!}\mathcal{O}_{U_i}$ corresponds to
$\sum_{k \in K'} s_k$. This finishes the proof.
\end{proof}

\end{document}
